% Generated draft: requires mathematical and source review.
\documentclass[11pt]{article}
\usepackage[margin=25mm]{geometry}
\usepackage{fontspec}
\setmainfont{Latin Modern Roman}
\usepackage{amsmath,amssymb,mathtools,unicode-math}
\setmathfont{Latin Modern Math}
\usepackage{enumitem,needspace,xurl,hyperref}
\setlength{\parindent}{0pt}
\setlength{\parskip}{5pt}
\setlength{\emergencystretch}{3em}
\newcommand{\N}{\mathbb{N}}
\newcommand{\Z}{\mathbb{Z}}
\newcommand{\Q}{\mathbb{Q}}
\newcommand{\R}{\mathbb{R}}
\newcommand{\C}{\mathbb{C}}
\newcommand{\F}{\mathbb{F}}
\newcommand{\A}{\mathbb{A}}
\newcommand{\PP}{\mathbb{P}}
\newcommand{\E}{\mathbb{E}}
\newcommand{\eps}{\varepsilon}
\newcommand{\dd}{\,\mathrm{d}}
\begin{document}
\section*{UnsolvedMath ID: 10400138}
\textbf{Source problem code:} catalogue entry
\subsection{Definitions and mathematical statement}
The QHI in the source are quantum hyperbolic invariants. S. Baseilhac and R. Benedetti ask to develop a four-dimensional theory of these invariants based on Turaev's shadow theory.

In that theory, a shadow is described by a suitably embedded simple two-dimensional polyhedron together with compatible gleams on its regions; gleams encode normal-bundle/thickening information needed to recover four-dimensional data. The exact shadow category and reconstruction hypotheses depend on the class of manifolds and any boundary or auxiliary structures under consideration. They are not specified in the excerpt.

The task is therefore to construct an extension of the source's QHI using such shadow data, with well-definedness under the relevant shadow equivalences and a stated relationship to the original quantum-hyperbolic construction. The source gives no state sum, target algebra, coefficient system, admissible four-manifolds, or required boundary and flat-bundle data. Nor does ``theory'' explicitly require a numerical invariant, a functorial topological quantum field theory, or a particular gluing law. These are unresolved design choices, not extra conditions supplied here. A generic four-dimensional quantum invariant without the QHI and shadow relationship would address a different task.
\subsection{Short English statement}
Develop a four-dimensional extension of the quantum hyperbolic invariants using Turaev shadows. The intended manifold category, shadow data, outputs, and relation to the original QHI are not specified.
\subsection{Sources}
\begin{itemize}
\item Local UnsolvedMath registry, permanent ID 10400138, source code catalogue entry; statement attribution S. Baseilhac and R. Benedetti; proposer field and proposal year absent; update 2026-08-06.
\item Ohtsuki, \emph{Problems on invariants of knots and 3-manifolds} (2002), Problem 7.23, PDF page 115: \url{https://msp.org/gtm/2002/04/gtm-2002-04-024s.pdf}. Registry access date 2026-07-29; PDF font-structure extraction. Provenance distinguishes 231 numbered bold headings. Present primary-literature status and redistribution terms require review.
\item Turaev's shadow theory is explicitly named, but no exact shadow reference, QHI definition, or later resolving paper is supplied. The research summary only repeats the programme.
\end{itemize}
\end{document}
